\section{RQ2 Customer Inactivity Pipeline}

The selected customer-snapshot export contains 72,146 observations for
10,613 customers. The supplied chronological split contains 50,982 TRAIN
rows from 31 December 2023 through 31 December 2024 and 21,164 TEST rows
at 31 March and 30 June 2025. The outcome is one when no order occurs in
the 90-day interval after the snapshot cutoff. It was reconstructed from
order headers and matched the supplied label on all 72,146 rows.

The feature policy uses only fields present in the historical snapshot and
deterministic transformations of those fields. Predictors include recency,
30-, 90- and 180-day order counts, 90-day sales and basket value, lifetime
orders, tenure, frequency and recency ratios, calendar seasonality and the
snapshot's home-shop identifier. Customer IDs, names, snapshot IDs, split
labels and future horizon fields are excluded. Current customer-master
attributes without effective-date history are also excluded because their
historical availability cannot be established.

Rolling-origin validation uses 30 June, 30 September and 31 December 2024.
The training window for each origin excludes labels whose 90-day outcome
window has not finished by that origin. A recency ranking rule, logistic
regression, LightGBM and CatBoost are compared. Platt calibration and the
F2 decision threshold are fitted from training-only out-of-time predictions.
The model is selected on validation PR-AUC, with calibrated Brier score as
the tie-breaker. The supplied TEST split is opened only after this selection.
The implementation writes aggregate metrics, date-level metrics,
calibration bins, customer-cluster bootstrap intervals and feature
importance; it does not write row-level customer predictions.

\section{RQ3 Replenishment Scenario Audit}

The supplied comparison table contains 2,500 paired cases across 32 shops
and 477 SKUs at a single cutoff date, 31 July 2025. Its future-demand field
is explicitly labelled synthetic oracle demand. The audit joins the product
master, verifies that both policy quantities are multiples of the case pack,
checks the supplied cost-gap identity, and calculates paired differences
in scenario lost units, average stock, supplied scenario cost and order
quantity. Uncertainty is estimated by resampling shops, preserving the
within-shop case structure.

The source does not document the scenario-cost formula and does not provide
the participating owner's actual current replenishment policy. Accordingly,
the experiment is a paired scenario audit rather than a historical replay
of the owner's policy. It cannot establish realised savings, causal policy
effects or the full RQ3 claim.

\section{RQ4 Decision Rule and Explanation Pipeline}

The richer RQ4 export contains 1,800 prepared cases across 32 shops and
467 SKUs, all dated 31 December 2025. The rule audit joins the product
master and independently recomputes inventory position as on-hand stock
plus pending purchases. Reordering is triggered when this position is below
the reorder point. The target is twice the reorder point; the required
quantity is rounded up to a case-pack multiple and capped by the number of
packs affordable under the supplied budget. Suggested quantity, action and
constraint reason are compared with the supplied fields.

A deterministic bilingual explanation generator has separate branches for
sufficient stock, a feasible reorder and a budget-blocked reorder. It states
the evidence relevant to each branch, including stock components, inventory
position, reorder point, target, pack arithmetic, feasible quantity, cost,
budget and supplier lead time. A numeric-coverage audit checks generated
English and Bangla explanations without exporting case-level text. The
constant human-rating field is not treated as participant evidence; the
counterbalanced owner study remains planned rather than completed.

\begin{figure}[H]
\centering
\begin{tikzpicture}[scale=0.86,transform shape,
source/.style={draw=dugreen,rounded corners,fill=softgreen,text width=4.2cm,minimum height=0.8cm,align=center,font=\footnotesize},
gate/.style={draw=dugold,rounded corners,fill=softorange,text width=4.2cm,minimum height=0.8cm,align=center,font=\footnotesize},
module/.style={draw=dugreen,rounded corners,fill=white,text width=4.4cm,minimum height=1.0cm,align=center,font=\footnotesize},
result/.style={draw=dugreen,rounded corners,fill=softblue,text width=5.2cm,minimum height=0.8cm,align=center,font=\footnotesize},
flow/.style={-{Latex[length=2mm]},thick,draw=gray!75}]
\node[source] (src) at (0,0) {Controlled SME data exports and dataset manifest};
\node[gate] (audit) at (0,-1.35) {Schema, provenance, join and point-in-time eligibility checks};
\node[module] (rq1) at (-3.7,-3.25) {RQ1: observed-sales forecasting\\LightGBM, CatBoost and statistical baselines};
\node[module] (rq2) at (3.7,-3.25) {RQ2: 90-day inactivity prediction\\recency, logistic regression, LightGBM and CatBoost};
\node[module] (rq3) at (-3.7,-5.45) {RQ3: paired replenishment scenarios\\cost, stock and lost-unit trade-offs};
\node[module] (rq4) at (3.7,-5.45) {RQ4: deterministic reorder rules\\bilingual evidence-based explanations};
\node[result] (out) at (0,-7.45) {Aggregate evidence, uncertainty and explicit promotion gates\\Human review remains responsible for business actions};
\draw[flow] (src)--(audit);
\draw[flow] (audit)--(rq1); \draw[flow] (audit)--(rq2);
\draw[flow] (audit)--(rq3); \draw[flow] (audit)--(rq4);
\draw[flow] (rq1)--(out); \draw[flow] (rq2)--(out);
\draw[flow] (rq3)--(out); \draw[flow] (rq4)--(out);
\end{tikzpicture}
\caption{Task-specific research pipelines under a shared evidence and data-quality gate.}
\label{fig:rq1_rq4_research_pipelines}
\end{figure}

The shared source package does not imply that every table is a valid
predictor or outcome for every research question. Each task retains its own
target definition, temporal boundary, evidence class and claim limit.
